<?php
// survey.php - Umfrageseite mit Validierung und Formular

// Datenbankverbindung einbinden
require __DIR__ . '/inc/db.php';

// Namen aus GET-Parameter auslesen
$name = isset($_GET['name']) ? trim($_GET['name']) : '';

// ============================================================
// 1. Fehlerbehandlung: Name existiert bereits in der DB
// ============================================================
$error = null;
$userId = null;
$displayName = '';

if ($name !== '') {
    // Prüfen, ob der Name bereits in user_table existiert
    $stmt = $pdo->prepare("SELECT Id FROM user_table WHERE Name = ?");
    $stmt->execute([$name]);
    $existingUser = $stmt->fetch();

    if ($existingUser) {
        $error = "Nutzer existiert schon, bitte anderen Namen wählen!";
    }
}

// ============================================================
// 2. Leerer Parameter -> anonymous
// ============================================================
if ($name === '' && $error === null) {
    // Prüfen, ob anonymous bereits existiert
    $stmt = $pdo->prepare("SELECT Id FROM user_table WHERE Name = 'anonymous'");
    $stmt->execute();
    $anonymousUser = $stmt->fetch();

    if ($anonymousUser) {
        // anonymous existiert bereits -> ID verwenden
        $userId = $anonymousUser['Id'];
        $displayName = 'anonym';
    } else {
        // anonymous existiert noch nicht -> neu anlegen
        $stmt = $pdo->prepare("INSERT INTO user_table (Name) VALUES ('anonymous')");
        $stmt->execute();
        $userId = $pdo->lastInsertId();
        $displayName = 'anonym';
    }
}

// ============================================================
// 3. Neuer Name (existiert noch nicht in der DB)
// ============================================================
if ($name !== '' && $error === null) {
    // Prüfen (sicherheitshalber nochmal), ob der Name existiert
    $stmt = $pdo->prepare("SELECT Id FROM user_table WHERE Name = ?");
    $stmt->execute([$name]);
    $existingUser = $stmt->fetch();

    if (!$existingUser) {
        // Neuen User anlegen
        $stmt = $pdo->prepare("INSERT INTO user_table (Name) VALUES (?)");
        $stmt->execute([$name]);
        $userId = $pdo->lastInsertId();
        $displayName = htmlspecialchars($name, ENT_QUOTES, 'UTF-8');
    } else {
        // sollte nicht passieren, da wir oben schon prüfen, aber zur Sicherheit
        $userId = $existingUser['Id'];
        $displayName = htmlspecialchars($name, ENT_QUOTES, 'UTF-8');
    }
}

// ============================================================
// 4. Wenn Fehler aufgetreten ist -> Meldung anzeigen
// ============================================================
if ($error !== null) {
    ?>
    <!DOCTYPE html>
    <html lang="de">
    <head>
        <meta charset="UTF-8">
        <meta name="viewport" content="width=device-width, initial-scale=1.0">
        <title>Fehler - Umfrage</title>
    </head>
    <body>
        <h1>Umfrage zum Studium</h1>
        <p style="color: red; font-weight: bold;"><?php echo htmlspecialchars($error, ENT_QUOTES, 'UTF-8'); ?></p>
        <p><a href="index.php">Zurück zur Startseite</a></p>
    </body>
    </html>
    <?php
    exit;
}